\documentclass[10pt,a4paper]{amsart}
\usepackage[margin=19mm]{geometry}
\usepackage{amsmath,amssymb,mathtools}
\usepackage[scr=rsfs,cal=euler]{mathalfa}
\usepackage{xfrac}
\usepackage[unicode,hidelinks,bookmarks=false]{hyperref}
\newcommand{\ie}{i.e.\ }
\newcommand{\cf}{cf.\ }
\newcommand{\resp}{respectively\ }
\newcommand{\Subsection}{\subsection}
\providecommand{\nobreakdash}{\nobreak\hbox{-}}
\setlength{\emergencystretch}{2em}
\multlinegap0pt
\usepackage{bm}
\usepackage{bbm}
\usepackage{dsfont}
\usepackage{xspace}
\usepackage{cancel}
\usepackage{mathtools}
\usepackage{amsthm}
\makeatletter
\@ifundefined{thm}{\newtheorem{thm}{Theorem}}{}
\@ifundefined{lem}{\newtheorem{lem}[thm]{Lemma}}{}
\makeatother
\title[Emergence for diffeomorphisms with nonzero Lyapunov exponent]{Emergence for diffeomorphisms with nonzero Lyapunov exponents}
\author{Agnieszka Zelerowicz}
\date{Source: arXiv:2111.07216v2 (CC BY 4.0)}
\begin{document}
\maketitle

\noindent\emph{Source and licence.} Emergence for diffeomorphisms with nonzero Lyapunov exponents. Version arXiv:2111.07216v2 (CC BY 4.0), distributed under the Creative Commons Attribution 4.0 International licence, as recorded in the arXiv licence metadata for this exact version and shown on the abstract page https://arxiv.org/abs/2111.07216v2. Full rights notice: licenses/arxiv-2111.07216-CC-BY-4.0.txt.\par\smallskip

\noindent\textit{Editorial scope.} Counted unit: the Lem whose source label is \texttt{lem:partsum} in the article (source0.tex lines 454--466, source label \texttt{lem:partsum}), delivered together with the article's own complete original proof of it (source0.tex lines 468--525, one \texttt{proof} environment of 58 lines). No mathematics is written, repaired or re-derived: statement and proof are the article's own text line for line. Required context retained uncounted: thm:LY (lines 173--180); thm:Pes (lines 385--416); eqn:norms (lines 392--401); eqn:temp (lines 392--401); thm:entr (lines 439--442). Every definition, display and label of those lines is retained unchanged. Omitted from this package and not redistributed: 1--172, 181--384, 402--438, 443--453, 467--467, 526--1430. Nothing inside a retained excerpt is omitted or altered: every retained body is delivered exactly as the article's lines are written, so no omission receipt of any kind accompanies this package. Citations: the retained excerpts carry no citation command at all, so no bibliography entry of the article is transcribed and no reference is redistributed. Reference adaptation: none was needed and none was made; every reference inside the delivered unit resolves inside it, so no unresolved reference or citation is produced. Printed slips kept literal: no printed word, symbol, hypothesis or number of the counted statement or of its proof was altered. Layout and class: the article's own class is \texttt{amsart} with packages including amsmath, amsthm, amscd, amssymb, mathrsfs, graphicx, enumerate, ascmac, color, here; the wrapper is the lane's pinned \texttt{amsart} editorial wrapper at 10pt on A4, which restates the theorem-like environments the retained text needs and adds the article's own macro definitions that the retained text needs (none), with extra wrapper packages (bm, bbm, dsfont, xspace, cancel, mathtools). The delivered numbering is the unit's own. Assets: no retained span contains a figure environment, a graphics-inclusion command, a PDF-page inclusion, a table or a TikZ picture; no image file is copied into this package, the package declares no asset dependency and no image file is redistributed with it. Licence: version arXiv:2111.07216v2 of this article is distributed under the Creative Commons Attribution 4.0 International licence, as recorded in the arXiv licence metadata for this exact version and shown on the abstract page https://arxiv.org/abs/2111.07216v2. A full rights notice is delivered at licenses/arxiv-2111.07216-CC-BY-4.0.txt. Characters: every retained excerpt is pure ASCII, so the delivered engine, XeLaTeX with its default Latin Modern text font, reports no missing character.
\begin{thm}\label{thm:LY}
There are numbers $d_1, d_2, \ldots, d_l$ such that:
\begin{enumerate}
\item $0\leq d_i\leq n_i$,
\item $\sum_{i=1}^l d_i = d_{\mu}^u$, and
\item $h_{\mu}(f)= \sum_{i=1}^l  d_i ~ \chi_i$.
\end{enumerate}
\end{thm}

\begin{thm}\label{thm:Pes}
Let $f : M \to M$ be a $C^{1+\alpha}$ diffeomorphism preserving an ergodic probability measure $\mu$
with Lyapunov exponents $\lambda_1 > \ldots > \lambda_s$ of multiplicities $k_1, \ldots k_s$.
Then there exists an invariant set $Z\subset M$ with $\mu(Z)=1$ such that for any $\eta>0$ the following holds on $Z$:
\begin{enumerate}
\item There exists a measurable family of invertible linear maps $C(x): T_xM \to \mathbb{R}^{\dim M}$ and $A_i(x)\in GL(k_i, \mathbb{R})$, $i=1, \ldots ,s,$ such that

\begin{enumerate}
\item\label{eqn:Oseledecsplit}
$$Df(x)= C^{-1}(f(x)) diag[ A_1(x), \ldots, A_s(x)  ] C(x),$$

\item\label{eqn:norms}
$$  e^{\lambda_i-\eta} < ||A_i^{-1}||^{-1} \leq ||A_i|| < e^{\lambda_i+\eta} ,  $$ 

\item\label{eqn:temp} 
$C(x)$ is tempered, that is for all $x\in Z$ $$  \lim_{m\to \pm \infty  }  \frac{1}{m}   \log ||   C(f^m(x))  ||   =    \lim_{m\to \pm \infty  }  \frac{1}{m}   \log ||   C^{-1}(f^m(x))  ||  =0;      $$
\end{enumerate}

\item there exist measurable functions $r,K: Z \to (0,1],$ and a collection of embeddings $\psi_x:  B(0, r(x)) \to M$ such that:

\begin{enumerate}
\item   $\psi_x(0) = x $ and $\psi_x = \exp_x \circ C(x),$

\item the maps $f_x = \psi_{f(x)}    \circ f   \circ \psi^{-1}_x$ satisfy   $d_{C^1}(f_x, Df_x(0))<\eta,$

\item there exists a constant $Q>0$ such that  for any $y,y'\in   \psi^{-1}_x(B(0,r(x))),$
$$      \frac{1}{K(x)}d(y,y')   \leq    ||  \psi_x(y) -  \psi_x(y')  ||   \leq    \frac{1}{Q}d(y,y')  ,$$

\item  $$  e^{-\eta}  < \frac{r(f(x))}{r(x)} < e^{\eta}  \text{ and } e^{-\eta}  < \frac{K(f(x))}{K(x)} < e^{\eta}  .$$
\end{enumerate}
\end{enumerate}
\end{thm}

\begin{thm}\label{thm:entr}
For any $\beta \in (0,1) ,$
$$  h_{\mu}(f) =      \lim_{\rho\to 0}    \lim_{n\to \infty}  \frac{1}{n} \log   C(n,\rho,\beta).$$
\end{thm}

\begin{lem}\label{lem:partsum}
If $\eta >0$ and $X\subset Z$ is a uniformity block of tolerance $\eta$, then there exists $m_0\in\mathbb{N}$ such that:
\begin{enumerate}
\item for all $m\geq m_0$ and any $x\in X$ one has\footnote{ We use the notation $A = \pm B$ to mean $  -B \leq  A  \leq B  $.} 
$$    \exp( S_m\phi (x)) =   | \det Df^m(x)_{| E^u(x)}  | ^{-d} =     e^{  \pm 3 \eta m d   \dim E^u}     \prod_{i=1}^{l}   e^{-\chi_i    d n_i   m };$$
\item for all $m\geq m_0$ and any finite set $A\subset X$,
$$  \frac{1}{m} \log  \sum_{x\in A}   \exp(S_m \phi (x)) > -3\eta d \dim(E^u) + \left[  \frac{1}{m} \log   \# A   - h_{\mu}(f)  \right] ; $$
\item for all $m\geq m_0$ and for any $\beta\in (0,1)$ one has
$$  \frac{1}{m} \log   Q(m,\rho,\beta, X) > -3\eta d \dim(E^u) + \left[  \frac{1}{m} \log   C(m,\rho,\beta)   - h_{\mu}(f)  \right] , $$
\item and in particular
$$  \lim_{\rho\to 0}   \liminf_{m\to \infty} \frac{1}{m} \log   Q(m,\rho,\beta, X) > -3\eta d \dim(E^u).$$
\end{enumerate}
\end{lem}

\begin{proof}
By Theorem \ref{thm:Pes}, for any $x\in Z$ we have that
$$  Df(x)_{| E^u(x)} = C^{-1}(f(x))  diag[   A_1(x), \ldots  ,  A_l(x)] C(x),  $$
consequently, since $Z$ is $f$-invariant,


\begin{align*}
  Df^m(x)_{| E^u(x)} & =    Df(f^{m-1}(x))   _{| E^u(f^{m-1}(x))}  \circ   \ldots   \circ    Df(f(x))   _{| E^u(f(x))}   \circ    Df(x) _{| E^u(x)}   \\
&= C^{-1}(f^m(x))  diag[  \prod_{k=1}^{m-1}  A_1(f^k(x)), \ldots  , \prod_{k=1}^{m-1}  A_l(f^k(x))] C(x).
\end{align*}



We then have that

$$   | \det Df^m(x)_{| E^u(x)}  | =  |\det C^{-1}(f^m(x)) | \cdot   \prod_{k=1}^{m-1} |\det  A_1(f^k(x))|\cdot \ldots  \cdot \prod_{k=1}^{m-1}  | \det A_l(f^k(x)) | \cdot |\det C(x)|. $$

We estimate all the terms in the product above.
By Statement (\ref{eqn:temp}) in Theorem \ref{thm:Pes} and because $X$ is a uniformity block, there exists $\tilde{m}_0\in\mathbb{N}$ such that
 $||   C( f^m(x))||, ||C^{-1}( f^m(x))||=  e^{ \pm \eta m}$ for all $m\geq \tilde{m}_0$. Consequently, by Hadamard's determinant inequality,
$$    |\det C^{-1}(f^m(x)) | =  e^{ \pm \eta m   \dim E^u}  \text{ for all $m\geq \tilde{m}_0$}. $$
By continuity of $C$ and $C^{-1}$ on $X$, there exists a constant $Q>0$ such that $  |\det C(x)|   = Q^{\pm 1}. $
In addition, for $i=1,\ldots, l$ and any $k\in\mathbb{N},$ Statement (\ref{eqn:norms}) of Theorem \ref{thm:Pes} implies that 
$$  |\det A_i(f^k(x))|   = e^{(\chi_i \pm \eta) n_i}. $$
Together we obtain that 


$$     | \det Df^m(x)_{| E^u(x)}  | =  e^{ \pm \eta m   \dim E^u}   \prod_{i=1}^{l}   e^{(\chi_i \pm \eta) n_i   m } Q^{\pm 1}   
 =    e^{ \pm 2 \eta m   \dim E^u} Q^{\pm 1}    \prod_{i=1}^{l}   e^{\chi_i  n_i   m }  $$
for every $x\in X$ and $m\geq \tilde{m}_0$. We set $m_0\geq \tilde{m}_0$ so that $Q< e^{\eta \dim E^u m_0}.$
Then for any $m\geq m_0$ one has
$$     \exp (S_m\phi (x)) =   | \det Df^m(x)_{| E^u(x)}  | ^{-d} =     e^{  \pm 3 \eta m d   \dim E^u}     \prod_{i=1}^{l}   e^{-\chi_i    d n_i   m } $$
so we proved the first statement of the lemma.
If now $    A \subset X$ is any finite set, then for all $m\geq m_0$ we can estimate

$$  \sum_{x\in A} \exp (S_m\phi (x)) = \sum_{x\in A }    | \det Df^m(x)_{| E^u(x)}  | ^{-d}   \geq     e^{ -  3 \eta m d   \dim E^u}    \prod_{i=1}^{l}   e^{-\chi_i    d n_i   m }  \# A.   $$

From this we obtain

$$   \frac{1}{m} \log   \sum_{x\in A} \exp( S_m\phi (x))    \geq    - 3 \eta  d   \dim E^u      - \sum_{i=1}^l  \chi_i    d n_i    +   \frac{1}{m} \log \# A$$
$$ =    - 3 \eta  d   \dim E^u    + \left[  - \sum_{i=1}^l  \chi_i    d n_i  + h_{\mu}(f)\right]   +  \left[ \frac{1}{m} \log  \# A  - h_{\mu}(f)\right]. $$

We complete the proof of the second statement of the lemma by observing that $   - \sum_{i=1}^l  \chi_i    d n_i  + h_{\mu}(f) = 0$ by Theorem \ref{thm:LY}.

%*****************************************8
If now $    E\subset X \text{ is }  \mu-(m,\rho,\beta)  \text{-spanning},$ then $\# E\geq   C(m,\rho,\beta) $ so Statement (3) is a consequence of Statement (2).
In addition, by Theorem \ref{thm:entr}, we conclude


\begin{align*}
 \lim_{\rho \to 0} \liminf_{m\to \infty}   \frac{1}{m} \log   Q(m,\rho,\beta, X)    &   \geq   \lim_{\rho\to 0}  \liminf_{m\to\infty}  \left[   - 3 \eta  d   \dim E^u- \sum_{i=1}^l  \chi_i    d n_i    +   \frac{1}{m} \log  C(m,\rho,\beta) \right] \\
 & = - 3 \eta  d   \dim E^u     - \sum_{i=1}^l  \chi_i    d n_i  + h_{\mu}(f) =  - 3 \eta  d   \dim E^u.
\end{align*}




\end{proof}

\end{document}
